\documentclass[landscape,a0b,final]{a0poster}
\usepackage[utf8]{inputenc}
\usepackage[ngerman]{babel}
\usepackage{amssymb}
\usepackage{graphicx}
\usepackage{hyperref}
\hypersetup
{
colorlinks=false,
linkcolor=false,
urlcolor=black,
breaklinks=true
}
\usepackage{tikz}
\usepackage{type1cm}
\usepgflibrary{arrows}
\usetikzlibrary{arrows}
\usetikzlibrary{shapes,snakes}
\usetikzlibrary[arrows]
\begin{document}
\begin{tikzpicture}[rounded corners][>=triangle 60]
\draw [step=1cm, white, very thin] (0,0) grid (114,84);
%\draw [step=1cm, gray, very thin] (0,0) grid (114,84);
%Auflistung der Institute
\filldraw [fill=white, draw=black, semitransparent] (1,0.5) rectangle (112,4.5);
\draw (56.5,2.5) node[text width=109cm, text justified]{\tiny 
\href{http://www.uni-luebeck.de/fakultaeten.php}{Institute der Universität zu
Lübeck}:\\ \href{http://www.anae.mu-luebeck.de}{Anaesthesiologie} \(|\) 
\href{http://www.anat.uni-luebeck.de/}{Anatomie} \(|\) 
\href{http://www.uk-sh.de/index.phtml?NavID=676.470.3&}{Arbeitsmedizin} \(|\)  \href{http://www.uk-sh.de/index.phtml?NavID=676.464.3&}{Augenheilkunde} \(|\) 
\href{http://www.biochem.uni-luebeck.de/}{Biochemie} \(|\) 
\href{http://www.bioweb.uni-luebeck.de/}{Biologie} \(|\) 
\href{http://www.bmo.uni-luebeck.de/}{Biomedizinische Optik} \(|\) 
\href{http://www.imbs-luebeck.de/}{Biometrie, Medizinische, und Statistik} \(|\) 
\href{http://www.chemie.uni-luebeck.de}{Chemie} \(|\) 
\href{http://www.chirurgie.uni-luebeck.de/}{Chirurgie} \(|\) 
\href{http://www.uk-sh.de/index.phtml?NavID=676.467.3&}{Dermatologie und Venerologie} \(|\) 
Frauenheilkunde und Geburtshilfe \(|\) 
\href{http://www.uk-sh.de/index.phtml?NavID=676.465.3&}{Hals-, Nasen- und Ohrenheilkunde mit Abteilung für Phoniatrie und Pädaudiologie} \(|\) 
\href{http://www.uk-sh.de/index.phtml?NavID=1595.502.3&}{Herzchirurgie} \(|\) 
\href{http://www.uk-sh.de/index.phtml?NavID=676.451.3&}{Humangenetik} \(|\) 
\href{http://www.fz-borstel.de/de/research/ibm/index.htm}{Abteilung Immunchemie und Biochemische Mikrobiologie am Forschungszentrum Borstel} \(|\)
\href{http://www.immunologie.uni-luebeck.de/}{Immunologie und Transfusionsmedizin} \(|\) 
\href{http://www.ifis.uni-luebeck.de/}{Informationssysteme} \(|\) 
\href{http://www.uk-sh.de/index.phtml?NavID=676.466.3&}{Kiefer- und Gesichtschirurgie} \(|\) 
\href{http://www.paedia.uni-luebeck.de/}{Kinder- und Jugendmedizin} \(|\) 
Kinder- und Jugendpsychiatrie \(|\) 
\href{http://www.uk-sh.de/?NavID=1595.503.3}{Kinderchirurgie} \(|\) 
\href{http://www.uk-sh.de/output/La3/676.1467.html}{Klinische Chemie - Zentrallabor} \(|\) 
\href{http://www.math.uni-luebeck.de/}{Mathematik} \(|\) 
\href{http://www.imgwf.uni-luebeck.de/_rubric/}{Medizin- und Wissenschaftsgeschichte} \(|\) 
\href{http://www.medinf.uni-luebeck.de/}{Medizinische Informatik} \(|\) 
\href{http://www.innere1.uni-luebeck.de/klinik/index.html}{Medizinische Klinik I} \(|\) 
\href{http://www.uk-sh.de/index.phtml?NavID=676.434.3&}{Medizinische Klinik II} \(|\) 
\href{http://www.uk-sh.de/index.phtml?NavID=676.435.3&}{Medizinische Klinik III} \(|\) 
\href{http://www.imt.uni-luebeck.de/}{Medizinische Mess- und Automatisierungstechnik (Medizintechnik)} \(|\) 
\href{http://www.uk-sh.de/index.phtml?NavID=676.460.3&}{Mikrobiologie, Medizinische, und Hygiene} \(|\) 
\href{http://www.molmed.uni-luebeck.de/}{Molekulare Medizin} \(|\) 
\href{http://www.imis.uni-luebeck.de/}{Multimediale und Interaktive Systeme} \(|\) 
\href{http://www.inb.uni-luebeck.de/}{Neuro- und Bioinformatik} \(|\) 
\href{http://www.uk-sh.de/index.phtml?NavID=676.443.3&}{Neurochirurgie} \(|\) 
\href{http://www.kfg.uni-luebeck.de/}{Neuroendokrinologie} \(|\) 
\href{http://www.neuro.uni-luebeck.de/}{Neurologie} \(|\) 
Neuroradiologie \(|\) 
\href{http://www.uk-sh.de/index.phtml?NavID=676.442.3&}{Orthopädie} \(|\) 
\href{http://www.patho.uni-luebeck.de/}{Pathologie} \(|\) 
\href{http://www.pharma.uni-luebeck.de/}{Pharmakologie und Toxikologie} \(|\) 
\href{http://www.physik.uni-luebeck.de/drupal/}{Physik} \(|\) 
\href{http://www.physio.uni-luebeck.de/}{Physiologie} \(|\) 
\href{http://www.uk-sh.de/index.phtml?mNavID=676.685&La=3}{Plastische Chirurgie} \(|\) 
\href{http://www.uk-sh.de/index.phtml?NavID=676.468.3&}{Psychiatrie und Psychotherapie} \(|\) 
\href{http://www.medpsy.uni-luebeck.de/}{Medizinische Psychologie} \(|\) 
\href{http://www.uk-sh.de/index.phtml?NavID=676.2388.3}{Radiologie und Nuklearmedizin} \(|\) 
\href{http://www.rmed.uni-luebeck.de/}{Rechtsmedizin} \(|\) 
\href{http://www.uk-sh.de/index.phtml?NavID=677.249.3&}{Rheumatologie} \(|\) 
\href{http://www.rob.uni-luebeck.de/}{Robotik und Kognitive Systeme} \(|\) 
\href{http://www.isip.uni-luebeck.de/}{Signalverarbeitung und Prozessrechentechnik} \(|\) 
\href{http://www.isp.uni-luebeck.de/}{Softwaretechnik und Programmiersprachen} \(|\) 
\href{http://www.uk-sh.de/index.phtml?NavID=676.471.3&}{Sozialmedizin} \(|\) 
\href{http://www.strahlentherapie-lübeck.de/}{Strahlentherapie} \(|\) 
\href{http://www.iti.uni-luebeck.de}{Technische Informatik} \(|\) 
\href{http://www.itm.uni-luebeck.de/}{Telematik} \(|\) 
\href{http://www.tcs.uni-luebeck.de/}{Theoretische Informatik} \(|\) 
\href{http://www.uk-sh.de/index.phtml?NavID=676.441.3&}{Unfallchirurgie} \(|\) 
\href{http://www.uk-sh.de/index.phtml?NavID=676.2390.3}{Urologie} \(|\) 
\href{http://www.vuz.uni-luebeck.de/}{Virologie und Zellbiologie}
};
%Mitglieder Kasten
\filldraw [fill=white, draw=black, semitransparent] (1,5) rectangle (85.5,12);
\draw (6.5,5.5) node[text width=10cm, text justified] {Mitglieder der Universität};
%Medizinische Fakultät
\filldraw [fill=orange!40!white, draw=orange!50!black, semitransparent] (2,6.5) rectangle (19,11.5);
\draw (7.5,7) node[text width=10cm, text justified] {\small Sektion Medizin};
\filldraw [fill=red!40!white, draw=red!50!black, semitransparent] (2.5,7.5) rectangle (7.5,11);
\draw (5,10.5) node[text width=4.5cm] {\small Professoren};
\filldraw [fill=red!40!white, draw=red!50!black, semitransparent] (8,7.5) rectangle (13,11);
\draw (10.5,9.5) node[text width=4.5cm, text justified] {\small Wiss. Mitarbeiter \& Assistenten};
\filldraw [fill=red!40!white, draw=red!50!black, semitransparent] (13.5,7.5) rectangle (18.5,11);
\draw (16,10) node[text width=4.5cm, text justified] {\small Weitere \\ Angestellte};
%Mitarbeiter
\filldraw [fill=gray!40!white, draw=gray!50!black, semitransparent] (20,6.5) rectangle (25,11.5);
\draw (22.5,9.5) node[text width=4.5cm, text justified] {\small Universitäts\\mitarbeiter außerhalb der Fakultät};
%TNF
\filldraw [fill=cyan!40!white, draw=cyan!50!black, semitransparent] (26,6.5) rectangle (42.5,11.5);
\draw (34,7) node[text width=15cm, text justified] {\small Sektion MINT};
\filldraw [fill=blue!40!white, draw=blue!50!black, semitransparent] (26.5,7.5) rectangle (31.5,11);
\draw (29,10.5) node[text width=4.5cm, text justified] {\small Professoren};
\filldraw [fill=blue!40!white, draw=blue!50!black, semitransparent] (32,7.5) rectangle (37,11);
\draw (34.5,9.5) node[text width=4.5cm, text justified] {\small Wiss. Mitarbeiter \& Assistenten};
\filldraw [fill=blue!40!white, draw=blue!50!black, semitransparent] (37.5,7.5) rectangle (42,11);
\draw (40,10) node[text width=4.5cm, text justified] {\small Weitere \\ Angestellte};
%Studenten
\filldraw [fill=white, draw=black, semitransparent] (43,5.5) rectangle (84.5,11.5);
\draw (48.5,6) node[text width=10cm, text justified] {Studenten};
%Graduiertenschule
\filldraw [fill=green!40!white, draw=green!50!black, semitransparent] (43.5,7.5) rectangle (49.5,11);
\draw (46.25,10) node[text width=5cm, text justified] {\small Doktoranten (GS)};
%TNF Studenten
%Fachschaft CS
\filldraw [fill=cyan!05!white, draw=cyan!50!black, semitransparent] (50.3,7.3)
rectangle (76.2,11.2);
%CLS
\filldraw [fill=blue!40!white, draw=blue!50!black, semitransparent] (50.5,7.5) rectangle (56.5,11);
\draw (53.25,10.5) node[text width=5cm, text justified] {\small CLS Studenten};
%Informatik
\filldraw [fill=blue!40!white, draw=blue!50!black, semitransparent] (57,7.5) rectangle (63,11);
\draw (59.75,10) node[text width=5cm, text justified] {\small Informatik\\Studenten};
%MIW
\filldraw [fill=blue!40!white, draw=blue!50!black, semitransparent] (63.5,7.5) rectangle (69.5,11);
\draw (66.25,10.5) node[text width=5cm, text justified] {\small MIW Studenten};
%Fachschaft MLS
%MLS
\filldraw [fill=blue!40!white, draw=blue!50!black, semitransparent] (70,7.5) rectangle (76,11);
\draw (73.25,10.5) node[text width=5cm, text justified] {\small MLS Studenten};
%Studenten MF
%Medizin
\filldraw [fill=red!40!white, draw=red!50!black, semitransparent] (77.5,7.5) rectangle (83.5,11);
\draw (80.25,10) node[text width=5cm, text justified] {\small Medizin\\Studenten};
%%ENDE MITGLIEDER
%Fachschaft Medizin
\filldraw [fill=orange!40!white, draw=orange!50!black, semitransparent] (72.5,15) rectangle (82.5,25);
\draw (76,16) node[text width=6cm, text justified]
{Fachschaftsvertretung\\Medizin}; \filldraw [fill=red!40!white, draw=red!50!black, semitransparent] (73,17) rectangle (77.5,20);
\draw (75.25,19) node[text width=4cm, text justified] {\small Gewählte\\Mitglieder};
\filldraw [fill=blue!10!white, draw=white!90!green, semitransparent] (77,19.5) circle (0.4);
\draw (77.3,19.5) node[text width=1cm, text justified] {\tiny 10};
\filldraw [fill=red!40!white, draw=red!50!black, semitransparent] (77.5,17) rectangle (82,20);
\draw (79.75,19) node[text width=4cm, text justified] {\small Freie\\Mitglieder};
\filldraw [fill=orange!40!white, draw=orange!50!black, semitransparent] (72.75,20.75) rectangle (82.25,24.75);
\draw (76,21.5) node[text width=6cm, text justified] {Vorsitz};
\filldraw [fill=red!40!white, draw=red!50!black, semitransparent] (73,23) rectangle (77,24.5);
\draw (77.75,24) node[text width=9cm, text justified] {\small Vorsitzender};
\filldraw [fill=red!40!white, draw=red!50!black, semitransparent] (77.5,23) rectangle (82,24.5);
\draw (82.25,24) node[text width=9cm, text justified] {\small Vize Vors.};
\filldraw [fill=red!40!white, draw=red!50!black, semitransparent] (77.5,21) rectangle (82,22.5);
\draw (82.23,22) node[text width=9cm, text justified] {\small Kassenwart};
\draw [->] (75.25,20) .. controls (75, 20.5) and (75.5, 20.5)..(75.25,20.75);
%Wähler Med
\draw [->] (80.5,11) .. controls (80.5, 13) and (76, 12) ..(77.5,15);
%Fachschaften der TNF
%\filldraw [fill=cyan!05!white, draw=cyan!50!black, semitransparent]
%(50.8,14.8) rectangle (71.7,25.2);
%Fachschaft CS
\filldraw [fill=cyan!40!white, draw=cyan!50!black, semitransparent] (51,15) rectangle (61,25);
\draw (55,16) node[text width=7cm, text justified]
{Fachschaftsvertretung \\ MINT}; \filldraw [fill=blue!40!white,
draw=blue!50!black, semitransparent] (51.5,17) rectangle (56,20); \draw (53.75,19) node[text width=4cm, text justified] {\small Gewählte\\Mitglieder};
\filldraw [fill=blue!10!white, draw=white!90!green, semitransparent] (55.7,19.7) circle (0.2);
\draw (56.1,19.7) node[text width=1cm, text justified] {\tiny 17};
\filldraw [fill=blue!40!white, draw=blue!50!black, semitransparent] (56,17) rectangle (60.5,20);
\draw (58.25,19) node[text width=4cm, text justified] {\small Freie\\Mitglieder};
\filldraw [fill=blue!40!white, draw=blue!50!black, semitransparent] (51.5,22.5) rectangle (60.5,24.5);
\draw (56.25,24) node[text width=9cm, text justified] {\small Sprecher};
\draw [->] (54.25,20) .. controls (54, 21) and (54.5, 21)..(54.25,22.5);
%Wähler CS
\draw [->] (54,11.2) .. controls (54, 13) and (57, 12) ..(56,15);
%StuPa
\filldraw [fill=yellow!40!white, draw=yellow!50!black, semitransparent] (62.5,43) rectangle (82.5,56);
\draw (70.25,43.5) node[text width=15cm, text justified] {Studierenden Parlament (StuPa)};
%Mitglieder
\filldraw [fill=yellow!70!white, draw=yellow!30!black, semitransparent] (63,44.5) rectangle (82,47);
\draw (70.75,45) node[text width=15cm, text justified] {Gewählte Mitglieder};
\filldraw [fill=blue!10!white, draw=white!90!green, semitransparent] (81.5,46.5) circle (0.4);
\draw (81.8,46.5) node[text width=1cm, text justified] {\tiny 23};
%Präsidium
\filldraw [fill=yellow!40!white, draw=yellow!50!black, semitransparent] (63.5,50) rectangle (81.5,55);
\draw (69,50.5) node[text width=10cm, text justified] {Präsidium};
\filldraw [fill=yellow!70!white, draw=yellow!30!black, semitransparent] (64.5,53) rectangle (72.25,54.5);
\draw (67,53.5) node[text width=5cm, text justified] {\small Präsident};
\filldraw [fill=yellow!70!white, draw=yellow!30!black, semitransparent] (72.75,53) rectangle (80.5,54.5);
\draw (75.25,53.5) node[text width=5cm, text justified] {\small Vize-Präsident};
\filldraw [fill=yellow!70!white, draw=yellow!30!black, semitransparent] (64,51) rectangle (69.75,52.5);
\draw (66.75,51.5) node[text width=5cm, text justified] {\small Gremienkoord.};
\filldraw [fill=yellow!70!white, draw=yellow!30!black, semitransparent] (70.25,51) rectangle (75.25,52.5);
\draw (73,51.5) node[text width=5cm, text justified] {\small Schriftführer};
\filldraw [fill=yellow!70!white, draw=yellow!30!black, semitransparent] (75.75,51) rectangle (81,52.5);
\draw (78.5,51.5) node[text width=5cm, text justified] {\small Wahlkoord.};
%Präsidiumswahl
\draw [->] (72.5,47) .. controls (72, 48) and (73, 49) ..(72.5,50);
%StuPaWahl
\draw [->] (82.75,11.5) .. controls (86, 16) and (85, 42) ..(80.5,43);
%Ausschusswahlen
%wahl in das Studentenwerk duch das StuPa
\draw [->, very thick] (71.25,56) to (71.25,58);
%AStA Wahl
\draw [->, very thick] (75,43) to (75,40);
%Asta Rekrutierung
\draw [->, very thick, dotted] (82,11.5) .. controls (83, 16) and (85, 20) ..(81,30);
%AStAs
\filldraw [fill=yellow!40!white, draw=yellow!50!black, semitransparent] (62.5,30) rectangle (82.5,40);
\draw (70.25,30.5) node[text width=15cm, text justified] {Allgemeiner Studierendenausschuss (AStA)};
%Gewählte Mitglieder
\filldraw [fill=yellow!40!white, draw=yellow!50!black, semitransparent] (63,31.5) rectangle (78.5,39.5);
\draw (67.25,32) node[text width=8cm, text justified] {\small Gewählte Mitglieder};
%Vorstand
\filldraw [fill=yellow!40!white, draw=yellow!50!black, semitransparent] (63.5,36.5) rectangle (75,39);
\draw (66.5,37) node[text width=5cm, text justified] {\small Vorstand};
\filldraw [fill=yellow!70!white, draw=yellow!30!black, semitransparent] (64,37.5) rectangle (69.5,38.5);
\draw (66.5,38) node[text width=5cm, text justified] {\small Vorsitzender};
\filldraw [fill=yellow!70!white, draw=yellow!30!black, semitransparent] (70,37.5) rectangle (74.5,38.5);
\draw (72.5,38) node[text width=5cm, text justified] {\small Vize-Vorsitz};
%Die Referate
\filldraw [fill=yellow!40!white, draw=yellow!50!black, semitransparent] (63.5,32.5) rectangle (66,36);
\draw (65,33.75) node[text width=3cm, text justified] {\small Ref.\\KOPR};
\filldraw [fill=yellow!40!white, draw=yellow!50!black, semitransparent] (66.5,32.5) rectangle (69,36);
\draw (68,33.75) node[text width=3cm, text justified] {\small Ref.\\KUS};
\filldraw [fill=yellow!40!white, draw=yellow!50!black, semitransparent] (69.5,32.5) rectangle (72,36);
\draw (71,33.75) node[text width=3cm, text justified] {\small Ref.\\EDV};
\filldraw [fill=yellow!40!white, draw=yellow!50!black, semitransparent] (72.5,32.5) rectangle (75,36);
\draw (74,33.75) node[text width=3cm, text justified] {\small Ref.\\PAS};
\filldraw [fill=yellow!40!white, draw=yellow!50!black, semitransparent] (75.5,32.5) rectangle (78,39);
\draw (77,33.75) node[text width=3cm, text justified] {\small Ref.\\Finanz};
%Sekreteriat
\filldraw [fill=yellow!70!white, draw=yellow!30!black, semitransparent] (79,30.5) rectangle (82,32);
\draw (80.75,30.75) node[text width=3cm, text justified] {\tiny Sekreteriat};
%Freie Mitarbeiter
\filldraw [fill=yellow!40!white, draw=yellow!50!black, semitransparent] (79,34.5) rectangle (82,39);
\draw (80.75,35.5) node[text width=3cm, text justified] {\tiny Freie\\Mitarbeiter};
%Wahl der Freien Mitarbeiter
\draw [->, very thick] (78.5,33) to (81,33) to (81,34.5);
%Konvent der Medizinischen Fakultät
\filldraw [fill=orange!40!white, draw=orange!50!black, semitransparent] (3,15) rectangle (13,25);
\draw (7.75,15.5) node[text width=9cm, text justified] {Senatsausschuss Med.};
%Dekanate
\filldraw [fill=red!40!white, draw=red!50!black, semitransparent] (3.5,23.5) rectangle (12.5,24.5);
\draw (5.5,24) node[text width=4cm, text justified] {Dekan};
\filldraw [fill=red!40!white, draw=red!50!black, semitransparent] (3.5,21.5) rectangle (7.5,22.5);
\draw (5.5,22) node[text width=4cm, text justified] {Prodekan};
\filldraw [fill=red!40!white, draw=red!50!black, semitransparent] (8.5,21.5) rectangle (12.5,22.5);
\draw (10.5,22) node[text width=4cm, text justified] {Prodekan};
\filldraw [fill=red!40!white, draw=red!50!black, semitransparent] (3.5,19.5) rectangle (12.5,20.5);
\draw (5.5,20) node[text width=4cm, text justified] {Studiendekan};
\filldraw [fill=red!40!white, draw=red!50!black, semitransparent] (3.5,16.5) rectangle (12.5,18.5);
\draw (5.5,17.5) node[text width=4cm, text justified] {Gewählte Mitglieder};
\filldraw [fill=blue!10!white, draw=white!90!green, semitransparent] (12.2,18.2) circle (0.2);
\draw (12.6,18.2) node[text width=1cm, text justified] {\tiny 4};
%Ausschusse des Medizinischen Konvents
\filldraw [fill=orange!40!white, draw=orange!50!black, semitransparent] (15,24) rectangle (25,25);
\draw (19.75,24.5) node[text width=9cm, text justified] {Promotionsausschuss};
\draw [->, very thick] (13,24.5) to (15,24.5);
\filldraw [fill=orange!40!white, draw=orange!50!black, semitransparent] (15,22.5) rectangle (25,23.5);
\draw (19.75,23) node[text width=9cm, text justified] {Forschungskommission};
\draw [->, very thick] (13,23) to (15,23);
\filldraw [fill=orange!40!white, draw=orange!50!black, semitransparent] (15,21) rectangle (25,22);
\draw (19.75,21.5) node[text width=9cm, text justified] {Strukturkommission};
\draw [->, very thick] (13,21.5) to (15,21.5);
%\filldraw [fill=orange!40!white, draw=orange!50!black, semitransparent] (15,22) rectangle (25,23);
%\draw (19.75,22.5) node[text width=9cm, text justified] {Auschuss 2};
%\draw [->, very thick] (13,22.5) to (15,22.5);
\filldraw [fill=orange!40!white, draw=orange!50!black, semitransparent] (15,18) rectangle (25,20);
\draw (19.75,18.5) node[text width=9cm, text justified] {Gleichstellungskomission};
\filldraw [fill=blue!10!white, draw=white!90!green, semitransparent] (24.7,19.7) circle (0.2);
\draw (25.1,19.7) node[text width=1cm, text justified] {\tiny 2};
\draw [->, very thick] (13,19) to (15,19);
\filldraw [fill=yellow!40!white, draw=yellow!50!black, semitransparent] (15,15) rectangle (24,17);
\draw (19.75,15.5) node[text width=9cm, text justified] {Habilitationsuasschuss};
\draw [->, very thick] (13,16) to (15,16);
\draw [->, very thick] (26,16) to (24,16);
%Wahl des Medizinischen Konvents
\draw [->] (8,11.5) .. controls (9, 13) and (7, 14) ..(8,15);
\draw [->] (78,11) .. controls (80, 17) and (6, 12) ..(8,15);
%Konvent TNF
\filldraw [fill=cyan!40!white, draw=cyan!50!black, semitransparent] (26,15) rectangle (36,25);
\draw (30.75,15.5) node[text width=9cm, text justified] {Senatsausschuss MINT};
%Dekanate
\filldraw [fill=blue!40!white, draw=blue!50!black, semitransparent] (26.5,23.5) rectangle (35.5,24.5);
\draw (28.5,24) node[text width=4cm, text justified] {Dekan};
\filldraw [fill=blue!40!white, draw=blue!50!black, semitransparent] (26.5,21.5) rectangle (30.5,22.5);
\draw (28.5,22) node[text width=4cm, text justified] {Prodekan};
\filldraw [fill=blue!40!white, draw=blue!50!black, semitransparent] (31.5,21.5) rectangle (35.5,22.5);
\draw (33.5,22) node[text width=4cm, text justified] {Prodekan};
\filldraw [fill=blue!40!white, draw=blue!50!black, semitransparent] (26.5,19.5) rectangle (35.5,20.5);
\draw (28.5,20) node[text width=4cm, text justified] {Studiendekan};
\filldraw [fill=blue!40!white, draw=blue!50!black, semitransparent] (26.5,16.5) rectangle (35.5,18.5);
\draw (28.5,17.5) node[text width=4cm, text justified] {Gewählte Mitglieder};
\filldraw [fill=blue!10!white, draw=white!90!green, semitransparent] (35.2,18.2) circle (0.2);
\draw (35.6,18.2) node[text width=1cm, text justified] {\tiny 4};
%ausschusse des Konvents der TNF
\filldraw [fill=cyan!40!white, draw=cyan!50!black, semitransparent] (38,24) rectangle (48,25);
\draw (42.75,24.5) node[text width=9cm, text justified] {Promotionsauschuss};
\draw [->, very thick] (36,24.5) to (38,24.5);
\filldraw [fill=cyan!40!white, draw=cyan!50!black, semitransparent] (38,22.5) rectangle (48,23.5);
\draw (42.75,23) node[text width=9cm, text justified] {Prüfungsuaschuss CLS};
\filldraw [fill=blue!10!white, draw=white!90!green, semitransparent] (47.7,23.2) circle (0.2);
\draw (48.1,23.2) node[text width=1cm, text justified] {\tiny 2};
\draw [->, very thick] (36,23) to (38,23);
\filldraw [fill=cyan!40!white, draw=cyan!50!black, semitransparent] (38,21) rectangle (48,22);
\draw (42.75,21.5) node[text width=9cm, text justified] {Prüfungsauschuss MLS};
\filldraw [fill=blue!10!white, draw=white!90!green, semitransparent] (47.7,21.7) circle (0.2);
\draw (48.1,21.7) node[text width=1cm, text justified] {\tiny 2};
\draw [->, very thick] (36,21.5) to (38,21.5);
\filldraw [fill=cyan!40!white, draw=cyan!50!black, semitransparent] (38,19.5) rectangle (48,20.5);
\draw (42.75,20) node[text width=9cm, text justified] {Prüfungsauschuss MIW};
\filldraw [fill=blue!10!white, draw=white!90!green, semitransparent] (47.7,20.2) circle (0.2);
\draw (48.1,20.2) node[text width=1cm, text justified] {\tiny 2};
\draw [->, very thick] (36,20) to (38,20);
\filldraw [fill=cyan!40!white, draw=cyan!50!black, semitransparent] (38,18) rectangle (48,19);
\draw (42.75,18.5) node[text width=9cm, text justified] {Prüfungsauschüsse Info.};
\filldraw [fill=blue!10!white, draw=white!90!green, semitransparent] (47.7,18.7) circle (0.2);
\draw (48.1,18.7) node[text width=1cm, text justified] {\tiny 6};
\draw [->, very thick] (36,18.5) to (38,18.5);
\filldraw [fill=cyan!40!white, draw=cyan!50!black, semitransparent] (38,16.5) rectangle (48,17.5);
\draw (42.75,17) node[text width=9cm, text justified] {Studienausschuss};
\filldraw [fill=blue!10!white, draw=white!90!green, semitransparent] (47.7,17.2) circle (0.2);
\draw (48.1,17.2) node[text width=1cm, text justified] {\tiny 2};
\draw [->, very thick] (36,17) to (38,17);
\filldraw [fill=cyan!40!white, draw=cyan!50!black, semitransparent] (38,15) rectangle (48,16);
\draw (42.75,15.5) node[text width=9cm, text justified] {Prüfungsaus. Grad. MLS};
\filldraw [fill=blue!10!white, draw=white!90!green, semitransparent] (47.7,15.7) circle (0.2);
\draw (48.1,15.7) node[text width=1cm, text justified] {\tiny 2};
\draw [->, very thick] (36,15.5) to (38,15.5);
%Gleichstellungskommision der TNF
\filldraw [fill=cyan!40!white, draw=cyan!50!black, semitransparent] (26,27) rectangle (36,29);
\draw (30.75,27.5) node[text width=9cm, text justified] {Gleichstellungskomission};
\filldraw [fill=blue!10!white, draw=white!90!green, semitransparent] (35.7,27.7) circle (0.2);
\draw (36.1,27.7) node[text width=1cm, text justified] {\tiny 1};
\draw [->, very thick] (31,25) to (31,27);
%Wahl zum Knovent TNF
\draw [->] (32,11.5) .. controls (33, 13) and (31, 14) ..(31,15);
\draw [->] (63.25,11.2) .. controls (63, 14) and (30, 11) ..(31,15);
%FUUU
%draw [->] (57,11.25) .. controls (57, 14) and (30, 11) ..(31,15);
%Senat
\filldraw [fill=yellow!40!white, draw=yellow!50!black, semitransparent] (03,43) rectangle (22.5,56);
\draw (7.75,43.5) node[text width=9cm, text justified] {Senat};
%Mitglieder
\filldraw [fill=yellow!70!white, draw=yellow!30!black, semitransparent] (3.5,44.5) rectangle (22,46.5);
\draw (11,45) node[text width=15cm, text justified] {Mitglieder};
\filldraw [fill=blue!10!white, draw=white!90!green, semitransparent] (21.7,46.2) circle (0.2);
\draw (22.1,46.2) node[text width=1cm, text justified] {\tiny 2};
%Senat nach Universitätsrat
\draw [->, very thick, dashed] (22.5,55) .. controls (28,56) and (40,55) ..(45,58);
%Ausschüsse des Senats
\filldraw [fill=yellow!40!white, draw=yellow!50!black, semitransparent] (4,27) rectangle (20,29);
\draw (9.75,27.5) node[text width=11cm, text justified] {\small Zentraler Gleichstellungsausschuss};
\filldraw [fill=yellow!40!white, draw=yellow!50!black, semitransparent] (4,29.5) rectangle (20,31.5);
\draw (8.75,30.5) node[text width=9cm, text justified] {\small Haushalts und Planungsausschuss};
\filldraw [fill=yellow!40!white, draw=yellow!50!black, semitransparent] (4,32) rectangle (20,34);
\draw (8.75,32.5) node[text width=9cm, text justified] {\small Zenraler Studienausschuss};
\filldraw [fill=yellow!40!white, draw=yellow!50!black, semitransparent] (4,34.5) rectangle (20,36.5);
\draw (8.75,35.5) node[text width=9cm, text justified] {\small Zentraler Ausschuss für Forschung und Wissenstransfer};
\filldraw [fill=yellow!40!white, draw=yellow!50!black, semitransparent] (4,37) rectangle (20,39);
\draw (8.75,38) node[text width=9cm, text justified] {\small Zentraler Haushalts- und Planungsausschuss};
\filldraw [fill=yellow!40!white, draw=yellow!50!black, semitransparent] (4,39.5) rectangle (20,41.5);
\draw (11.75,40.5) node[text width=15cm, text justified] {\small Ständige
Komission für die Sicherung der guten \\ wissenschaftlichen Praxis}; \draw
(8.75,38) node[text width=9cm, text justified] {\small Zentraler Haushalts- und Planungsausschuss}; \filldraw [fill=yellow!40!white, draw=yellow!50!black, semitransparent] (5,57) rectangle (20,59);
\draw (10.5,58) node[text width=10cm, text justified] {\small Koordinierungsausschuss Forschungszentrum Borstel};
\filldraw [fill=yellow!40!white, draw=yellow!50!black, semitransparent] (5,59.5) rectangle (20,61.5);
\draw (10.5,60.5) node[text width=10cm, text justified] {\small Koordinierungsausschuss Zentrale Hochschulbib.};
\filldraw [fill=yellow!40!white, draw=yellow!50!black, semitransparent] (5,62) rectangle (20,64);
\draw (10.5,63) node[text width=10cm, text justified] {\small Koordinierungsausschuss Studiengang Medizintechnik an der FHL};
%Wahl des Senats
\draw [->] (02,12) .. controls (01,14) and (01,40) ..(04,43);
\draw[-, very thick] (22,43) to (22,26);
\draw[-, very thick] (28,26) to (8,26);
\draw[->, very thick] (28,26) to (28,25);
\draw[->, very thick] (8,26) to (8,25);
\draw[->, very thick] (22,28) to (20,28);
\draw[->, very thick] (22,30.5) to (20,30.5);
\draw[->, very thick] (22,33) to (20,33);
\draw[->, very thick] (22,35.5) to (20,35.5);
\draw[->, very thick] (22,38) to (20,38);
\draw[->, very thick] (22,40.5) to (20,40.5);
\draw[-, very thick] (22,56) to (22,63);
\draw[->, very thick] (22,58) to (20,58);
\draw[->, very thick] (22,60.5) to (20,60.5);
\draw[->, very thick] (22,63) to (20,63);
%bild inner mitte
\draw (84.5,42) node[text width=109cm, text justified]
{
%\begin{figure}[htb]
%\centering
\includegraphics[scale=2.6]{nurlogo.jpg}
%\end{figure}
};
%Präsidium
\filldraw [fill=yellow!40!white, draw=yellow!50!black, semitransparent] (03,66) rectangle (22.5,73);
\draw (7.75,66.5) node[text width=9cm, text justified] {Präsidium};
%wahl des Präsidiums
\draw [->, very thick] (04,56) to (04,66);
\filldraw [fill=yellow!40!white, draw=yellow!50!black, semitransparent] (04,70) rectangle (12,72);
\draw (8.75,70.5) node[text width=9cm, text justified] {Kanzler};
\filldraw [fill=yellow!40!white, draw=yellow!50!black, semitransparent] (13.5,70) rectangle (21.5,72);
\draw (18.75,70.5) node[text width=9cm, text justified] {Präsident};
\filldraw [fill=yellow!40!white, draw=yellow!50!black, semitransparent] (13.5,67) rectangle (21.5,69);
\draw (18.75,67.5) node[text width=9cm, text justified] {Vize Präsident};
%Universitätsrat
\filldraw [fill=magenta!20!white, draw=magenta!70!black, semitransparent] (43,58) rectangle (62.5,65);
\draw (49.75,58.75) node[text width=13cm, text justified] {Universitätsrat Schleswig Holstein};
\filldraw [fill=magenta!40!white, draw=magenta!50!black, semitransparent] (43.5,60) rectangle (62,62);
\draw (48.25,60.5) node[text width=9cm, text justified] {Mitglieder};
\filldraw [fill=blue!10!white, draw=white!90!green, semitransparent] (61.7,61.7) circle (0.2);
\draw (62.1,61.7) node[text width=1cm, text justified] {\tiny 1};
%Medizin-Ausschuss
\filldraw [fill=magenta!20!white, draw=magenta!70!black, semitransparent] (33,58) rectangle (42.5,65);
\draw (38,58.75) node[text width=9cm, text justified] {Medizin-Ausschuss};
\filldraw [fill=magenta!40!white, draw=magenta!50!black, semitransparent] (33.5,60) rectangle (42,62);
\draw (38.25,60.5) node[text width=9cm, text justified] {Mitglieder};
%Studentenwerk
\filldraw [fill=magenta!20!white, draw=magenta!70!black, semitransparent] (63,58) rectangle (82.5,65);
\draw (70.75,58.75) node[text width=15cm, text justified] {Studentenwerk Schleswig Holstein};
\filldraw [fill=magenta!40!white, draw=magenta!50!black, semitransparent] (63.5,60) rectangle (82,62);
\draw (68.25,60.5) node[text width=9cm, text justified] {Verwaltungsrat};
\filldraw [fill=blue!10!white, draw=white!90!green, semitransparent] (81.7,61.7) circle (0.2);
\draw (82.1,61.7) node[text width=1cm, text justified] {\tiny 2};
\filldraw [fill=magenta!40!white, draw=magenta!50!black, semitransparent] (63.5,63.5) rectangle (72.5,64.5);
\draw (68.25,64) node[text width=9cm, text justified] {Vorstand};
\filldraw [fill=blue!10!white, draw=white!90!green, semitransparent] (72.2,64.2) circle (0.2);
\draw (72.6,64.2) node[text width=1cm, text justified] {\tiny 1};
\draw [->] (69,62) .. controls (68, 62.5) and (70, 62.6) ..(69,63.5);
\filldraw [fill=magenta!40!white, draw=magenta!50!black, semitransparent] (73,63.5) rectangle (82,64.5);
\draw (77.75,64) node[text width=9cm, text justified] {Geschäftsführer};
%Zusätzliche Einrichtungen
%Schülerakademie
\filldraw [fill=green!20!white, draw=green!70!black, semitransparent] (3,78) rectangle (12.5,81);
\draw (7.75,78.5) node[text width=9cm, text justified] {Schülerakademie};
%IT-Service Center
\filldraw [fill=green!20!white, draw=green!70!black, semitransparent] (3,74) rectangle (12.5,77);
\draw (7.75,74.5) node[text width=9cm, text justified] {IT-Service Center};
%Graduate School
\filldraw [fill=green!20!white, draw=green!70!black, semitransparent] (13,78) rectangle (22.5,81);
\draw (17.75,78.5) node[text width=9cm, text justified] {Graduate School};
%Fernuni
\filldraw [fill=green!20!white, draw=green!70!black, semitransparent] (13,74) rectangle (22.5,77);
\draw (17.75,75) node[text width=9cm, text justified] {Zentrum für Fernstudium und Weiterbildung};
%Alumni
\filldraw [fill=green!20!white, draw=green!70!black, semitransparent] (23,78) rectangle (32.5,81);
\draw (27.75,78.5) node[text width=9cm, text justified] {Alumni};
%Mensa
\filldraw [fill=green!20!white, draw=green!70!black, semitransparent] (23,74) rectangle (32.5,77);
\draw (27.75,74.5) node[text width=9cm, text justified] {Career Center};
%Zentrale Hochschulbibliotek
\filldraw [fill=green!20!white, draw=green!70!black, semitransparent] (23,70) rectangle (32.5,73);
\draw (27.75,71) node[text width=9cm, text justified] {Zentrale Hochschulbibliothek};
%Hochschulsport
\filldraw [fill=green!20!white, draw=green!70!black, semitransparent] (23,66) rectangle (32.5,69);
\draw (27.75,66.5) node[text width=9cm, text justified] {Hochschulsport};
%Gemeinsame Tierhaltung
\filldraw [fill=green!20!white, draw=green!70!black, semitransparent] (23,62) rectangle (32.5,65);
\draw (27.75,62.5) node[text width=9cm, text justified] {Gemeinsame Tierhaltung};
%Careeer Center
\filldraw [fill=green!20!white, draw=green!70!black, semitransparent] (23,58) rectangle (32.5,61);
\draw (27.75,58.5) node[text width=9cm, text justified] {Mensa};
%Universitätsklinikum
\filldraw [fill=gray!10!white, draw=gray!50!black, semitransparent] (33,70) rectangle (62.5,81);
\draw (41,70.75) node[text width=15cm, text justified] {Universitätsklinikum Schleswig Holstein};
\filldraw [fill=gray!40!white, draw=gray!50!black, semitransparent] (33.5,78.5) rectangle (42.5,80);
\draw (39.25,79) node[text width=11cm, text justified] {Aufsichtsrat};
\filldraw [fill=gray!40!white, draw=gray!50!black, semitransparent] (33.5,76.5) rectangle (42.5,78);
\draw (39.25,77) node[text width=11cm, text justified] {Vorstand};
\filldraw [fill=gray!40!white, draw=gray!50!black, semitransparent] (33.5,74) rectangle (42.5,75.5);
\draw (39.25,74.5) node[text width=11cm, text justified] {Personalräte};
\filldraw [fill=gray!40!white, draw=gray!50!black, semitransparent] (33.5,71.5) rectangle (42.5,73.5);
\draw (39.25,72) node[text width=11cm, text justified] {Mitarbeiter};
\filldraw [fill=gray!10!white, draw=gray!50!black, semitransparent] (43,71.5) rectangle (62,80);
\draw (48.25,72) node[text width=10cm, text justified] {Tochtergesellschaften};
\filldraw [fill=gray!40!white, draw=gray!50!black, semitransparent] (43.5,77.5) rectangle (61.5,79.5);
\draw (52,78) node[text width=16cm, text justified] {Bildungszentrum für Gesundheitsfachberufe};
\filldraw [fill=gray!40!white, draw=gray!50!black, semitransparent] (43.5,75) rectangle (61.5,77);
\draw (52,75.5) node[text width=16cm, text justified] {Medizinisches Versorgungszentrum};
\filldraw [fill=gray!40!white, draw=gray!50!black, semitransparent] (43.5,72.5) rectangle (61.5,74.5);
\draw (52,73) node[text width=16cm, text justified] {Zentrum für Integrative Psychiatrie};
%Landesregierung
\filldraw [fill=gray!10!white, draw=gray!50!black, semitransparent] (63,66) rectangle (82.5,83.5);
\draw (68,66.75) node[text width=9cm, text justified] {Landesregierung SH};
\filldraw [fill=gray!40!white, draw=gray!50!black, semitransparent] (63.5,68) rectangle (75.5,70);
\draw (69.5,69) node[text width=11cm, text justified] {Ministerium für Wissenschaft, Wirtschaft und Verkehr};
\filldraw [fill=gray!40!white, draw=gray!50!black, semitransparent] (63.5,70.5) rectangle (75.5,72.5);
\draw (69.5,71.25) node[text width=11cm, text justified] {Finanzministerium};
\filldraw [fill=gray!40!white, draw=gray!50!black, semitransparent] (63.5,73) rectangle (75.5,75);
\draw (69.5,73.75) node[text width=11cm, text justified] {Innenministerium};
\filldraw [fill=gray!40!white, draw=gray!50!black, semitransparent] (63.5,75.5) rectangle (75.5,77.5);
\draw (69.5,76.5) node[text width=11cm, text justified] {Ministerium für Bildung und Frauen};
\filldraw [fill=gray!40!white, draw=gray!50!black, semitransparent] (63.5,78) rectangle (75.5,80);
\draw (69.5,79) node[text width=11cm, text justified] {Ministerium für Justiz, Arbeit und Europa};
\filldraw [fill=gray!40!white, draw=gray!50!black, semitransparent] (63.5,80.5) rectangle (82,82.5);
\draw (72.5,81.5) node[text width=17cm, text justified] {Ministerium für Soziales, Gesundheit, Familie, Jugend und Senioren };
\filldraw [fill=gray!40!white, draw=gray!50!black, semitransparent] (76,68) rectangle (82,80);
\draw (78.25,69.5) node[text width=4cm, text justified] {Minister-präsident};
\draw [->, very thick] (63.5,69) to (60,65);
%%SIDEBAR
\filldraw [fill=white, draw=black, semitransparent] (86,81.5) rectangle (112,83.5);
\draw (99.5,82.5) node[text width=24.5cm, text justified] {\tiny Angefertigt von Lukas Ruge Stand 07.2009. Als Quelle dienen UniVis, Ordnungen und Gesetze sowie Informationen aus den Gremien. No rights reserved what so ever.};
%Relevantes aus den Ordungen
\filldraw [fill=white, draw=black, semitransparent] (86,33) rectangle (112,81);
\draw (98.5,57.9) node[text width=25.7cm, text justified] {
\tiny \begin{itemize}
\item  \textbf{\href{http://www.unirat-sh.uni-kiel.de/}{Universitätsrat}:} ist
ein gemeinsamer Hochschulrat der Universitäten Lübeck, Kiel und Flensburg. Er
besteht aus 9 Stimmberechtigten Mittgliedern. Im Universitätsrat sitzt
beratend auch einer (oder mehrere) studentische Vertreter der Uni Lübeck
(Üblicherweise der AStA Vorsitzende). Der Universitätsrat ist gleichzeitig der
Hochschulrat der Universitäteten und ist in seinen Aufgaben durch § 19 des HSG
sowie durch seine eigene Geschäftsordnung bestimmt.
\begin{itemize}
\item \textbf{\href{http://www.medizin-ausschuss.uni-kiel.de/}{Medizin-Ausschuss}:} 
Die Universitäten Kiel und Lübeck betreiben einen gemeinsamen
Medizinausschuss, dieser koordiniert das Leerangebot, entscheided über
Profesurauschreibungen, regelt die Finanzen und ist dem Universitätsrat über
selbige Verteilung jährlich einen Bericht schuldig. Er setzt sich zusammen aus
einem Wissenschaftsdirektor als Vorsitzendem, sowie Presidiumsvertretern der
Universitäten in Kiel und Lübeck, sowie de Vorstandsvorsitzenden des UKSH.
\end{itemize}
\item
\textbf{\href{http://www.uni-kiel.de/stwsh/seiten_allgemein/sitemap.html}{Studentenwerk
Schleswig Holstein}:} Ein Dienstleistungsunternehmen mit wirtschaftlicher,
sozialer und kultureller Zielsetzung eingerichtet nach dem Studentenwerkgesetz
von 1971. Es untersteht der Rechtsaufsicht des Ministeriums für Bildung,
Wissenschaft, Forschung und Kultur des Landes Schleswig-Holstein. Das
Studentenwerk betreibt unter anderem die Mensen und viele Wohnheime in
Schleswig Holstein. Zudem stellt es u.a. Bafög Beratung und psychische
Studienberatung zur Verfügung. Es finanziert sich u.a. durch die
Semestergebüren der Studenten. Das Magazin \textit{POTTKIEKER} wird vom
Studentewerk heraufsgebracht.
\begin{itemize}
\item \textbf{Der Verwaltungsrat:} Der Verwaltungsrat entscheidet in
gesetzlich bestimmten Angelegenheiten und wählt den Vorsitzenden und zwei
Studierende in den Vorstand. In den Verwaltungsrat entsendet jede Hochschule
je angefangener 3.000 der ihr angehörender Studierender zwei Vertreter/innen.
\item \textbf{Vorstand:} Der Vorstand leitet die Verwaltung des
Studentenwerks. Der Vorstand setzt sich zusammen aus dem
Vorstandsvorsitzenden, dem Geschäftsführer, einem studentischen Mitglied für
die wissenschaftlichen Hochschulen und einem studentischen Mitglied für die
Fachhochschulen und die Musikhochschule Lübeck.
\item \textbf{Geschäftführer:} Der Geschäftsführer vertritt das Studentenwerk
gerichtlich und rechtsgeschäftlich nach außen.
\end{itemize}
\item \textbf{\href{http://univis.uni-luebeck.de/prg?search=departments&name=praesidium&show=single}{Präsidium}:} (ehemals Rektorrat) Das Präsidium der Universität zu Lübeck besteht aus dem Präsidenten, dem Vizepräsidenten und dem Kanzler. Es wird nach der XXX Ordnung auf xxx weise gebildet / gewählt / bestätigt.
\item \textbf{Senat:} Der Senat besteht aus 13 Personen (7 Professoren, 2
Studenten, 2 Wis. Mitarbeitern, 2 n. Wiss Mitarb.) Seine
Entscheidungskompetenz liegt in der Wahl des Präsidiums, dem Treffen des
Haushaltsplanes der Herausgabe von Stellungnahmen, der Einrichtung von
Studienplätzen und dem Beschluss von Ordnungen. Der Senat vertritt nach außen
die gesamte Universität. Er setzt zudem die Senatsausschüsse der Sektionen ein.
\begin{itemize}
\item \textbf{Zentraler Gleichberechtigungsausschuss:} Stellt die
Gleichberechtigung an der Universität sicher.
\item \textbf{Haushalts und Planungsauschuss:} Unter Vorsitz des Kanzlers wird
hier der Haushalt der Universität erarbeitet und dem Senat vorgelegt.
\item \textbf{Senatsausschuss MINT:} Der Senatsauschuss seitzt sich aus dem
Dekan, 11 Professoren, 4 wissenschaftlichen Mitarbeitern, 2 Studenten (mit je
einem Vertreter) und 2 Mitarbeitern des nichtwissenschaftlichen Dienstes
zusammen. Es erarbeitet die Studiengänge an der Sektion und kümmert sich um
deren fachliche Belange. Gewählt wird der Ausschuss durch die Mitglieder der
Sektion MINT, allerdings besetzt duch den Senat, welcher die Wahl dabei
respektieren soll.
\begin{itemize}
\item \textbf{Prüfungsauschüsse:} Die Prüfungsausschüsse der Sektion erarbeiten
unter anderem die Prüfungsordnung für das jeweilige Studienfach.
\item \textbf{Gleichstellungskomision:} Zuständig für Überwachung der
Gleichstellung an der Universität
\item \textbf{Strukturkomision}
\end{itemize}
\item \textbf{Senatsausschuss Medizin:} Zusammensetzung und Aufgaben wie bei dem
Sektion MINT.
\begin{itemize}
\item \textbf{Promotionsauschuss}
\item \textbf{Forschungskommission}
\item \textbf{Strukturkommission} \\ (Die Strukturkomission der Sektion MINT ist
aus Platzgründen nicht abgebildet.)
\end{itemize}
\end{itemize}
\item \textbf{Gemeinsamer Habilitationsausschuss der Sektionen:} Berät über
Habilitationen.
\item \textbf{Das Studierendenparlament:} ist neben dem AStA eines der
zentralen Organe der Studierendenschaft. Es wird jährlich gewählt und dient
der zentralen Meinungs- und Willensbildung. Es Wählt die Mitglieder des AStA
und führt die Finanzielle Aufsicht über das Budget der Studierendenschaft. Es
beschließt und erlässt zudem die Ordnungen der Studierendenschaft und richtet
Fachschaften ein. Zudem hat das StuPa laut Vollvbersammlungsordung (VVO § 1)
die Möglichkeit Vollversammlungen einzuberufen. Das StuPa hat die Kompetenz
Ausschüsse zu bilden. Das StuPa tagt Monatlich in Vorlesungszeit.
\item \textbf{\href{http://www.asta.uni-luebeck.de/}{Allegeine
Studierendenauschuss}:} Der AStA führt die laufenden geschäfte der
Studierendenschaft und vertritt sie nach außen. Er hällt sich dabei an die
Maßgabe des StuPa. Der AStA tagt regelmäßig in der Vorlesungszeit. Der AStA
veröffentlicht das \textit{StudentenPack}.
\begin{itemize}
\item \textbf{Vorstand:} Der Vorsitzende des AStA ist Leiter des AStA und der
Vertreter der Studierendenschaft. Er wird durch das StuPa direkt gewählt. Der
Vorsitzende hat immer einen stellvertretenden Vorsitzenden.
\item \textbf{Referat Finanzen:} Als einziges Referat ist das Finanzreferat
zwangsläufig zu besetzen. Der Referatsleiter ist Vertreter des
Vize-Vorsitzenden.
\item \textbf{Die Referatsstruktur:} Der AStA teilt sich in Referate auf, die
jeweils eine Referatsleitung haben. Die Referate kümmern sich vorrangig um
bestimmte Bereiche der Belange der Studenten
\end{itemize}
\item \textbf{\href{http://medizin.fachschaften.uni-luebeck.de/}{Fachschaft
Medizin}:} Die Fachschaft Medizin vertritt die fachlichen Belange der
Medizin-Studenten der Universität zu Lübeck. Sie wird ein mal pro Jahr gewählt.
\item \textbf{\href{http://cs.fachschaften.uni-luebeck.de}{Fachschaft MINT}:}
Die Fachschaft MINT vertritt die fachlichen Belange der
Studenten der Sektion MINT in den Studienfächer MIW, MML, MLS Medizinische
Informatik und Informatik der Universität zu Lübeck. Sie wird ein mal pro Jahr
gewählt.
\item \textbf{Mitglieder der Universität}: Studenten, Angestellte, Dozenten,
Wissenschaftliche- und nichtwissenschaftliche Mitarbeiter sowie die Mitglieder
des Präsidiums sind Mitglieder der Universität.
\end{itemize}
\begin{itemize}
\item \textit{\href{http://www.itsc.uni-luebeck.de/}{IT-Service Center}}
Stellt Dienstleistungen für den Betrieb rechnergestützten
Informationsverarbeitung und Kommunikation für Forschung, Lehre, Studium und
Verwaltung der Universität. Das enthällt u.a. das Uni-LAN, die Rechner und die
Emails. Geleitet wird es von einer IT-Kommission, die keine studentischen
Mitglieder hat und nichtöffentlich tagt.
\item \textit{\href{http://www.gradschool.uni-luebeck.de/}{Graduate School for
Computing in Medicine and Life Sciences}} Ziel ist eine interdisziplinäre
Ausbildung von Doktoranden an den Schnittstellen Informatik und Medizin,
Informatik und Lebenswissenschaften sowie Lebenswissenschaften und Medizin.
Die Graduiertenschule wird von einem akademischen Leiter und einem
Geschäftsführer geführt und nach außen repräsentiert. Fragen, die die
Graduiertenschule als Gesamtheit betreffen, werden durch den Lenkungsausschuss
entschieden. Dieser setzt sich aus Vertretern der Lehrkräfte der Schule, der
Studierenden und der Universität zusammen. Die Studierenden werden durch zwei
Studentische Vertreter der Gradschool vertreten.
\item \textit{\href{http://www.zhb.uni-luebeck.de/}{Zentrale
Hochschulbibliothek}}: Die Zentrale Hochschulbibliothek Lübeck dient als
wissenschaftliche Fachbibliothek der Forschung, Lehre, dem Studium und der
beruflichen Weiterbildung auf dem Gebiete der Medizin, Naturwissenschaft und
der Technik und ist in erster Linie für die Literaturversorgung der
Hochschulangehörigen zuständig. Darüber hinaus steht die ZHB allen Einwohnern
Schleswig-Holsteins zur Nutzung offen.
\begin{itemize}
\item Für die Koordination der Bibliothek gibt es einen Runden Tisch an dem
Studenten, Mitarbeiter und Dozenten sitzen.
\end{itemize}
\item Mensa: Die Mensa der Universität zu Lübeck wird vom
Studentenwerk-Schleswig Holstein verwaltet.
\end{itemize}
};
%Legende
\filldraw [fill=white, draw=black, semitransparent] (86,5) rectangle (112,32);
\draw [->] (86.5,30) .. controls (87, 29) and (87, 31) ..(88,30);
\draw [->, very thick] (86.5,26.25) to (88,26.25);
\draw [->, very thick, dotted] (86.5,23) .. controls (87, 22) and (87, 24) ..(88,23);
\draw [->, very thick, dashed] (86.5,19) .. controls (87, 20) and (87, 18) ..(88,19);
\filldraw [fill=cyan!40!white, draw=cyan!50!black, semitransparent] (86.5,16) rectangle (88,17);
\filldraw [fill=orange!40!white, draw=orange!50!black, semitransparent] (86.5,13.9) rectangle (88,14.9);
\filldraw [fill=yellow!40!white, draw=yellow!50!black, semitransparent] (86.5,11.8) rectangle (88,12.8);
\filldraw [fill=green!20!white, draw=green!70!black, semitransparent] (86.5,9.7) rectangle (88,10.7);
\filldraw [fill=magenta!20!white, draw=magenta!70!black, semitransparent] (86.5,7.6) rectangle (88,8.6);
\filldraw [fill=blue!10!white, draw=white!90!green, semitransparent] (87.25,5.75) circle (0.2);
\draw (87.65,5.75) node[text width=1cm, text justified] {\tiny 4};
\draw (100.5,18.5) node[text width=20cm, text justified] {
\tiny
\textbf{Wahl nach § 17 HSG}: "`[...] werden die Mitglieder der Organe und sonstigen Gremien von den jeweiligen Mitgliedergruppen in freier, gleicher und geheimer Wahl unmittelbar gewählt."' soll heißen: Studenten wählen die Studentischen Vertreter, Professoren die Professoren usw. \\
\\
\textbf{Bestimmt}: Das Gremien wählt oder bestimmt die Mitglieder, allerdings nicht (zwangsläufig) aus seiner Mitte. Die Wahl des Präsidiums der Universität ist gesondert in der Präsidiumswahlordnung bestimmt.\\
\\
\textbf{Rekrutiert sich aus}: Sofern die Wählenden nicht auch die Gruppe aller mit passiven Wahrecht sind, ist bei einem Gremium angegeben aus welcher Personengruppe sich die Mitglieder zusammensetzen können \\
\\
\textit{Schlägt vor}: Durch dieses Gremum werden Personen Vorgeschlagen, welche dann von anderer Seite Bestimmt werden. \\
\\ 
\textbf{Der TNF Zugehörend} \\
\\
\textbf{Der Medizinischen Fakultät zugehörend} \\
\\
\textbf{Universitäre und Studentische Gremien \& Personengruppen außerhalb der Fakultäten}\\
\\
\textbf{Weitere Universitäre Einrichtungen} \\
\\
\textbf{Überuniversitäre Gremien und Organisatuonen sowie Landes- und Bundesorganisationen} \\
\textbf{Anzahl Studentischer Mitglieder}
};
\end{tikzpicture}
\newpage
Personenverzeichniss: \\
\begin{tabular}{|p{15cm}||p{15cm}||p{15cm}||p{15cm}||p{15cm}|||p{15cm}||p{15cm}|}
\hline
\multicolumn{3}{|c||}{Studentische Gremien} & \multicolumn{2}{||c|}{Konvente} & \multicolumn{2}{||c|}{Universitäre Gremien} \\
\hline
\hline
\textbf{Studierendenparlament} & \textbf{AStA} & \textbf{Fachschaften} & \textbf{Konvent der Med. Fak} &  \textbf{Konvent der TNF} & \textbf{Senat} & \textbf{Präsidium}\\
\begin{itemize}
\item Christoph Zabel
\item Alma Osmanovic
\item Georg Engelbart
\item Florian Markowsky
\item Nils Uflacker
\item Thiemo Sprink
\item Johanna Neuhold
\item Niklas Finck
\item Katharina Ilm
\item Frauke Riechert
\item Florian Hinze
\item Dennis Boldt
\item Christoph Leschcyk
\item Julius Naujoks
\item Teresa Pliet
\item Linda Krause
\item Julien Beck
\item Anna-Katharina Trull
\item Thierry Panje
\item Inga Stolz
\item Christoph Curths
\item Hannes Nauhaus
\item Sven-Thomas Antoni
\end{itemize}
&
\begin{itemize}
\item Vorsitz
\begin{itemize}
\item
\item
\end{itemize}
\item Finanzen
\begin{itemize}
\item
\item
\end{itemize}
\item Politik, Ausland, Soziales
\begin{itemize}
\item
\item
\end{itemize}
\item Kultur, Umwelt, Sport
\begin{itemize}
\item
\item
\end{itemize}
\item Elekronische Datenverarbeitung
\begin{itemize}
\item
\item
\end{itemize}
\item Öffentlichkeitsarbeit
\begin{itemize}
\item
\item
\end{itemize}
\item Freie Mitarbeiter
\begin{itemize}
\item
\item
\end{itemize}
\end{itemize}
&
\textbf{Fachschaft Medizin}
\begin{itemize}
\item Christoph Zabel (Sprecher)
\item Alma Osmanovic
\item Lena Kertelge
\item Verena Becker
\item Kim-Eva Linhoff
\item Bjorn Heuwer
\item Anna Beilharz
\item Mirja Muller
\item Christopth Marquetand
\item Felicitas Kahlitz
\end{itemize}
\textbf{Fachschaft MLS}
\begin{itemize}
\item Thiemos Sprink (Sprecher)
\item Julius Naujoks
\item Daniel Hasche
\item Patrick Beckert
\item Frauke Kausch
\end{itemize}
\textbf{fachschaftCS}
\begin{itemize}
\item Viola Borchard (Sprecher)
\item Anne Reichart
\item Jens Hocke
\item Anna Barkentien
\item Christian Klante
\end{itemize}
&
\begin{itemize}
\item
\end{itemize}
&
\begin{itemize}
\item 
\end{itemize}
&
\begin{itemize}
\item Prof. Dr. rer. nat. Walter Dosch
\item Prof. Dr. rer. nat. Stefan Fischer
\item Prof. Dr. rer. nat. Enno Hartmann
\item Prof. Dr. med. Hendrik Lehnert
\item Prof. Dr.-Ing. Erik Maehle, Erik
\item Prof. Dr. med. Werner Solbach
\item Prof. Dr. med. Jürgen Westermann 
\item apl. Prof. Dr. Stefan Anemüller
\item Dr. med. Hauke Nielsen 
\item \textit{Michael Drefahl}
\item \textit{Philipp Wewering}
\item \textit{Sebastian Wien} 
\item	Rena Giese
\item Dipl.-Ing. Axel Schiller
\end{itemize}
&
\begin{itemize}
\item Prof. Dr. med. Peter Dominiak, (Präsident)
\item Prof. Dr. rer. nat. Thomas Martinetz (VP) 
\item Dr. Oliver Grundei (Kanzler)
\end{itemize}\\
\hline
\hline
\end{tabular} \\
Weitere Studentische Vertreter (2008/2009): 
\begin{itemize}
	\item Universitäsrat Schleswig Holstein: Philip Wewering
	\item Auschüsse des Senats
	\begin{itemize} 
		\item Zentraler Haushalts- und Planungsausschuss: Martin Demmert und Arnika Wagner
		\item Zentraler Studienausschuss: ??
		\item Zentraler Ausschuss für Forschung und Wissenstransfer: Claudia Korn, Miriam Skilandat
		\item Zentraler Haushalts und Planungsausschuss: ??
		\item Zentraler Gleichstellungsausschuss: Katharina Moser-Notborn (\textit{ist das ein Senatsuasschuss})
		\item Kommission zur Förderung des wissenschaftlichen Nachwuchses: Michael Seydel
		\item Koordinierungsausschuss Forschungszentrum Borstel: 
	\end{itemize}
	\item Komisionen uns Auschüsse der TNF
	\begin{itemize}
		\item Gleichstellungskommission der TNF: Anja Kerstein
		\item Strukturkommission TNF: Thiemo Sprink, Julia Kinzel
		\item Promotionsauschuss der TNF: Phillip Wewering
		\item Prüfungsauschuss CLS: Anna Barlentien
		\item Prüfungsauschuss MLS: A. Barkentien und Johannes Lotz
		\item Prüfungsausschuss Graduiertenschule MLS: J. Naujoks und I. Choschzick
		\item Prüfungsauschuss MIW: J. Kuczewski  und T. Fritz
		\item \textbf{Prüfungsauschüsse} Informatik 
		\begin{itemize}
			\item Diplom: Jan Kolbaum, Thomas Klähn
			\item Bachelor: Christian Wolters
			\item Master: Jan Kolbaum, Thomas Klähn
		\end{itemize}
		\item Studienausschuss TNF: Anna Barkentien und Thiemo Sprink
		\item 
	\end{itemize}
	\item Prüfungsauschuss Medizin:
	\item Verterter im Studentenwerk Schleswig Holstein: Linda Krause
\end{itemize}
Inofizielle Posten:
\begin{itemize}
\item Kontaktperson für den Stadtverkehr:
\end{itemize}
\end{document}